\section{Results}
\subsection{Step 1: The Published Analysis Is Reproducible, but In-Sample Fit Is Not a Reliable Guide}\label{sec:repro}
\textbf{The data are exactly the data of the paper.} Table~\ref{tab:datacheck} lists every printed quantity that depends on the sample sizes or on individual values. The twelve printed BIC values are recovered exactly when the public numbers of data points $n_T$ (36, 28, 120, 64, 144 and 141) are inserted into the printed $F_\mathrm{obj}$ and $k$; a single missing or duplicated row would change them. The four fingerprint sums of squares (74.2, 840.6, 535.3 and 1616.0) are recovered to the printed decimal, which requires every individual value to match. The composition of the data sets (four studies, two temperatures, three water activities, three of five batches used) also agrees. All \NCheckPass{}/\NCheckAll{} checks passed. Consequently, every comparison in this paper is made on exactly the data the authors used, and any difference between methods is due to the method.
\inctab{datacheck}

\textbf{Step 1 of the original method is reproduced exactly.} Table~\ref{tab:reprofp} compares the printed and reproduced fingerprint estimates: all four pairs of selectivities and all four sums of squares agree to the printed precision, and the 95\,\% confidence intervals agree within 2--12\,\%, a difference attributable to the finite-difference Hessian used by the authors. Figure~\ref{fig:repro}a shows the same agreement graphically: the printed and reproduced values of $S_M$, $S_D$ and SS (circles) and of the simplified dynamic fits of Åkerman and Bornadel (squares) lie on the identity line over six orders of magnitude. The fingerprint estimates that serve as prior centres in PPK-FIP are therefore the authors' own numbers.
\inctab{repro_fp}

\textbf{Step 2 is reproduced where the printed fits are at their optimum, and reveals two problems where they are not.} Table~\ref{tab:reprodyn} gives, for each of the twelve dynamic fits, the printed $F_\mathrm{obj}$, its value when the printed parameters are inserted into our implementation, the implied enzyme loading $E=k_E/k_A$, and our global optimum. The simplified models of Åkerman et al. and the Bornadel models are reproduced to the last digit ($F_\mathrm{obj}$, BIC, $F_{AZ}$, $F_{AM}$, $F_{AD}$), and evaluating the model at the printed parameters recovers the printed $F_\mathrm{obj}$ in \ReproConsistent{} of \ReproAll{} tables. The implied loadings are round numbers, 2.00\,\% for Åkerman, 5.00\,\% for Bornadel and 15.0\,\% for Mao (Figure~\ref{fig:repro}c, where the points fall on the grid lines at 2, 5 and 15), which confirms that the parameterisation is identical; for the Mao case the printed full model is recovered ($F_\mathrm{obj}=\MaoFullAtPaper$ vs.\ 1475.9) only with the initial acid:TMP ratio of the actual charge, 10:3.3 ($A_0=303.0$), not with the nominal 3:1. Two discrepancies remain, and both are visible in Figure~\ref{fig:repro}b, which shows for every fit the printed objective (orange), the objective at the printed parameters (yellow) and our global optimum (blue). (i) Several printed full models are not optimal: the blue bar lies below the orange one for the Åkerman 100\,°C, Tao and Mao full models, most strikingly for Åkerman 100\,°C (\AkHundredFullGlobal{} vs.\ 117.9, a factor of 2.4). (ii) For the Tao case the yellow bars rise far above the orange ones: the printed dynamic fits (121.2/123.5) could not be reproduced from the published profiles, the best value at the printed parameters over all initial acid ratios being 131.1.
\inctab{repro_dyn}
\begin{figure}[!htbp]\centering\incfig{fig03_reproduction}\caption{Reproduction audit. (a) Printed versus reproduced fingerprint and simplified-model quantities (log--log). (b) $F_\mathrm{obj}$ printed, evaluated at the printed parameters, and at our global optimum (F: full model, S: simplified model). (c) Enzyme loading implied by $k_E/k_A$ at the printed parameters; grid lines at 0.1, 2, 5 and 15.}\label{fig:repro}\end{figure}

Figure~\ref{fig:landscape} explains why the full model is prone to non-optimal fits. Panel (a) shows the local optima reached from 200 random starts of the original full model for each case, as $\log_{10}$ of the ratio to the global minimum, together with the value printed in the paper (diamond). Only \LandMinFrac--\LandMaxFrac\,\% of the starts reach the global minimum, the rest end on plateaus up to three orders of magnitude higher, and the printed Åkerman 100\,°C value sits on one of the secondary optima ($\log_{10}$ ratio 0.38). Panel (b) shows the profile objective along $F_A$ with all other parameters re-optimised: for the three cases shown it decreases monotonically towards $F_A\to\infty$ and never turns up. The full model is therefore structurally non-identifiable \cite{raue2009,transtrum2015}: in the limit $F_A\to\infty$ the rate becomes independent of acid saturation and $k_E$ and the $F$ constants can only be estimated as ratios, which is also why the printed confidence intervals of the full models are so large and why our own full-set fits sit at the parameter bound (Table~\ref{tab:params} below).
\begin{figure}[!htbp]\centering\incfig{fig04_landscape_identifiability}\caption{(a) Local optima reached from 200 random starts of the original full dynamic model (dots), with the value printed in the paper (diamonds); numbers above give the percentage of starts that reach the global minimum. (b) Profile objective along $F_A$ (other parameters re-optimised), relative to its minimum: the objective decreases monotonically towards $F_A\to\infty$, i.e.\ the full model is not identifiable.}\label{fig:landscape}\end{figure}

\storylink{Two conclusions carry forward. First, the data and the reference implementation are exact, so every later difference between methods is attributable to the estimator. Second, the model-selection evidence of the original paper rests on in-sample statistics evaluated partly at non-optimal points of a non-identifiable model, and one case cannot be reproduced at all. This motivates three decisions: all comparisons use our global re-fits ORIG and ORIG-full as references, which are at least as good as the printed fit of the same model in every case; BIC is also reported with all estimated parameters counted; and the decisive criteria are out-of-sample (Step~3).}

\subsection{Step 2: The Remaining Misfit Has a Mechanistic Signature}
Table~\ref{tab:insample} and Figure~\ref{fig:insample} compare the fits to the full data, the protocol of the original paper. Figure~\ref{fig:insample}a normalises each $F_\mathrm{obj}$ by the best value printed in the paper for that case: the ORIG refit (orange) equals or exceeds 1 because the simplified form is more restrictive than the best printed form, the ORIG-full refit (yellow) is at or below 1 because it reaches the global optimum, and PPK-FIP (blue) is lowest in every case, at 0.96, 0.27, 0.32, 0.62, 0.80 and 0.79 of the printed value. Figure~\ref{fig:insample}b shows $\mathrm{BIC}_\mathrm{tot}$ relative to the re-fitted ORIG model, with all estimated parameters counted: PPK-FIP is lower in \InBICWins{} of 6 cases, by 14 to 123 units, and differs only for Åkerman 70\,°C (+6), where the penalty for the two mechanisms that the data set to zero is counted. The gains are large but unevenly distributed, and the rest of this section locates them.
\inctab{insample}
\begin{figure}[!htbp]\centering\incfig{fig08_insample}\caption{In-sample comparison. (a) $F_\mathrm{obj}$ relative to the best value printed in the paper. (b) BIC$_\mathrm{tot}$ (all estimated parameters counted) relative to the re-fitted original simplified model; negative is better.}\label{fig:insample}\end{figure}

Figure~\ref{fig:fits} shows the fitted curves of the re-fitted original method (dashed) and PPK-FIP (solid) for all nine profiles, Figure~\ref{fig:paths} the composition paths against $c$ for the three profiles in which the new mechanisms are active, Figure~\ref{fig:resid} the residuals, and Table~\ref{tab:params} the estimates.

\textbf{Åkerman et al.} At 70\,°C all models describe the nine sampling times well ($F_\mathrm{obj}$ \InAkSeventyFnew{} vs.\ 59.0 printed) and the two curves in Figure~\ref{fig:fits} are indistinguishable; PPK-FIP sets processivity to zero ($p_M=0$) and keeps deactivation at the prior level ($k_d=\AkSeventyKd$\,h$^{-1}$), so that its only difference from the original model is the BIC penalty for two mechanisms it correctly leaves inactive, a first illustration that PPK-FIP does not add mechanisms where the data do not call for them. At 100\,°C the improvement is large, from 117.9 to \InAkHundredFnew{} (\InAkHundredRed\,\%), and it comes from the second half of the profile, where the reaction slows down more than acid depletion alone can explain. PPK-FIP attributes this to first-order deactivation with $k_d=\AkHundredKd$\,h$^{-1}$ (BIC-weighted over the two model forms; 0.032\,h$^{-1}$ in the dominant simplified form, Table~\ref{tab:params}), i.e.\ a half-life of the active enzyme of roughly one day at 100\,°C, against no detectable loss at 70\,°C.

\textbf{Bornadel et al.} This is the case the original authors flagged as mass-transfer limited, and the data select pseudo-processivity decisively: $p_M=\BornPM$ with a 95\,\% half-width of 0.025 (Table~\ref{tab:params}), $p_D=\BornPD$, and $F_\mathrm{obj}$ falls from 329.9 to \InBornFnew{} (\InBornRed\,\%) while $\mathrm{BIC}_\mathrm{tot}$ falls by more than 120 units. Figure~\ref{fig:resid}b shows the mechanism at work in the residuals of M and D: the original model (circles, dashed) over-predicts M by 2--5 molar\,\% and under-predicts D by 4--5 molar\,\% during the first two hours, then reverses sign around 5\,h, a systematic pattern that no choice of lumped constants can remove because it is fixed by the composition path; with pseudo-processivity (squares, solid) the early residuals shrink to within $\pm2$ molar\,\% and lose their systematic sign. Figure~\ref{fig:paths} (left) shows the same thing in the $c$-domain: the fingerprint path of the original method (dashed) rises too early in M and too late in D, whereas the processive path (solid) follows the data. The selectivities change accordingly, from the apparent values $S_M=0.909$, $S_D=0.066$ to $S_M=\BornSM$, $S_D=\BornSD$, in agreement with the processive fingerprint analysis of Mitchell et al. \cite{mitchell2023}; in Figure~\ref{fig:fits} the two curves look alike at full scale because the difference is concentrated in the first eight hours and is spread over 30 sampling times.

\textbf{Tao et al.} At $a_w=0.25$, $F_\mathrm{obj}$ decreases from 121.2 (printed) to \InTaoAFnew{} (\InTaoARed\,\%), mainly through the full set of lumped constants (weight of the simplified form 0.02, Table~\ref{tab:params}) and a small processivity of the diester step ($p_D=\TaoAPD$); the middle panel of Figure~\ref{fig:paths} shows a correspondingly small shift of the path. At $a_w=0.35/0.45$ the decrease is from 325.8 to \InTaoBFnew{} (\InTaoBRed\,\%).

\textbf{Mao et al.} The three reuse batches are fitted with batch-specific $k_E$ as in the original paper; $F_\mathrm{obj}$ decreases from 1475.9 to \InMaoFnew{} (\InMaoRed\,\%). The remaining misfit is dominated by the late T and D values of batch~1 (Figure~\ref{fig:fits}, bottom left): the measured triester plateaus near 90\,\% while D stays at 5--10\,\% after 70\,h, whereas any model that conserves TMP and lets the reaction run to completion must send T to 100\,\% and D to zero; the residual that remains there is therefore a property of the digitised data rather than of the kinetics, as anticipated in Section~\ref{sec:methods}, and the same feature explains the interpolation results of Step~6. Figure~\ref{fig:resid}a summarises the residuals of all six cases: the median absolute residual of PPK-FIP is below that of the original model in every case, the largest reductions of the upper quartile are for Åkerman 100\,°C (2.7 to 1.1 molar\,\%) and Bornadel (1.7 to 1.0), and the Mao whiskers, which are set by the batch-1 plateau, shorten from 7.7 to 5.9.
\inctab{params}
\begin{figure}[!htbp]\centering\incfig{fig06_fits_time}\caption{Fits to the full data: original two-step method re-fitted to its global optimum (dashed) and PPK-FIP (solid). Markers: data.}\label{fig:fits}\end{figure}
\begin{figure}[!htbp]\centering\incfig{fig07_paths_c}\caption{Composition paths against the degree of reaction: fingerprint path of the original method (dashed) and PPK-FIP path with pseudo-processivity (solid); estimated processivities in the panel titles. Colours as in Figure~\ref{fig:fits}.}\label{fig:paths}\end{figure}
\begin{figure}[!htbp]\centering\incfig{fig09_residuals}\caption{(a) Absolute residuals of the full-data fits. (b) Residuals of M and D for the Bornadel profile: the systematic early pattern of the original model (circles, dashed) is removed by pseudo-processivity (squares, solid).}\label{fig:resid}\end{figure}

Table~\ref{tab:params} makes the pattern explicit across all profiles. The processivity $p_M$ is zero within its confidence half-width everywhere except Bornadel, where it is seven standard errors from zero; $p_D$ is small (0.01--0.06) and clearly non-zero only for Tao $a_w=0.25$ and Åkerman 100\,°C; $k_d$ stays at the prior level ($10^{-4}$--$10^{-3}$\,h$^{-1}$) everywhere except Åkerman 100\,°C; the selectivities move away from the fingerprint values only where processivity is present; and the full lumped-constant set, where it dominates, sits on the identifiability ridge ($k_E$ at its bound), confirming Figure~\ref{fig:landscape}b.

\storylink{The in-sample gains are large, but more importantly they are localised and interpretable. One mechanism, pseudo-processivity, acts in the mass-transfer-limited Bornadel data, exactly where the original authors anticipated it, and removes a residual pattern that lumped constants cannot touch; the other, deactivation, acts in the single profile that slows down late; and almost nothing changes where the original model was adequate (Åkerman 70\,°C). PPK-FIP therefore does not fit better everywhere through added flexibility; it fits better where a specific mechanism is missing. Whether this translates into prediction is the question of Step~3.}

\subsection{Step 3: The Gain Survives Out-of-Sample}
Lower in-sample error is expected for a richer model; the decisive question is whether it predicts data not used for fitting. Table~\ref{tab:loto} and Figure~\ref{fig:loto} report the leave-one-time-point-out cross-validation, in which both steps of the original method are re-estimated in every one of the \NLotoFolds{} folds. Figure~\ref{fig:loto}a shows the RMSE per case for the original simplified and full models, the joint fit without prior (JOINT-BMA) and PPK-FIP; Figure~\ref{fig:loto}b shows the ratio of PPK-FIP to each original form on a log axis, with Holm-corrected significance stars coloured by the direction of the paired median. Pooled over all \NFitPoints{} held-out observations, PPK-FIP reduced the RMSE from \LotoAllOrig{} to \LotoAllNew{} molar\,\% (\LotoAllRedOrig\,\%, $p$\,\LotoAllPOrig) relative to the original simplified method and from \LotoAllFull{} by \LotoAllRedFull\,\% ($p$\,\LotoAllPFull) relative to the original full model. Per case, PPK-FIP had a lower RMSE than the original method in \LotoWinsOrig{} of 6 cases (\LotoSigWinsOrig{} significant after Holm correction, none significantly worse) and lower than the full model in \LotoWinsFull{} of 6 (\LotoSigWinsFull{} significantly better). The largest gains are for Bornadel (\LotoBornOrig$\to$\LotoBornNew), Tao $a_w$ 0.25 (\LotoTaoAOrig$\to$\LotoTaoANew) and Mao (\LotoMaoOrig$\to$\LotoMaoNew), i.e.\ the cases in which Step~2 located a missing mechanism or a dominant full $F$ set; for Åkerman 70\,°C, the profile with only nine sampling times, the three estimators lie within 0.25 molar\,\% of one another (\LotoAkSeventyNew, \LotoAkSeventyOrig{} and \LotoAkSeventyFull). Figure~\ref{fig:loto}a also previews Step~4: the joint fit without prior is as good as PPK-FIP on the densely sampled profiles but clearly worse on the two sparse Åkerman profiles (2.34 vs.\ 1.89 and 3.88 vs.\ 3.51). A more conservative test that treats each fold rather than each observation as the unit confirms the result: PPK-FIP has the lower fold error in \FoldWinsOrig{} of \FoldN{} folds against the original method ($p=\FoldPOrig$) and in \FoldWinsFull{} of \FoldN{} against the full model ($p=\FoldPFull$).
\inctab{loto_main}
\begin{figure}[!htbp]\centering\incfig{fig10_loto}\caption{Leave-one-time-point-out cross-validation. (a) RMSE per case. (b) RMSE ratio PPK-FIP/reference; stars mark Holm-corrected Wilcoxon tests (* $p<0.05$, ** $p<0.01$, *** $p<0.001$), blue when the paired median favours PPK-FIP and orange otherwise.}\label{fig:loto}\end{figure}

Forward extrapolation probes behaviour beyond the training window (Table~\ref{tab:fcst}, Figure~\ref{fig:fcstbars}): trained on the first half of each profile, every method predicts the second half with errors of 1--20 molar\,\% instead of the 1--4 of interpolation. PPK-FIP halved the extrapolation error for Åkerman 100\,°C (\FcstAkHundredOrig$\to$\FcstAkHundredNew), where the deactivation identified in Step~2 is already partly expressed within the training window, reduced it for Tao $a_w$ 0.25 (\FcstTaoAOrig$\to$\FcstTaoANew) and for Bornadel and Tao $a_w$ 0.35/0.45, and pooled over all cases the two mechanistic estimators are equivalent (\FcstAllNew{} vs.\ \FcstAllOrig, $p=\FcstAllPOrig$). For Åkerman 70\,°C (\FcstAkSeventyOrig{} vs.\ \FcstAkSeventyNew) and the Mao batches (\FcstMaoOrig{} vs.\ \FcstMaoNew) the ranking reverses, and the reason is informative: the first half of a short profile contains the rise of M and D but not the late behaviour from which $k_d$ and the processivities are identified, so the protocol defines the data requirement of the additional mechanisms, namely a profile that extends into the slow-down, as the ablation of Step~4 confirms. The original full model illustrates the opposite problem, with results ranging from 0.67 (Åkerman 100\,°C) to 9.5 (Åkerman 70\,°C), as expected of a non-identifiable model whose ridge is pinned down differently by different halves of the data. Figure~\ref{fig:fcst} shows three representative profiles with the training window shaded. For Åkerman 100\,°C PPK-FIP (solid) comes closest to the measured T at 24\,h among the mechanistic models; for Bornadel both mechanistic models track T and D closely, whereas the neural ODE (dotted) lets M rise again after 15\,h, which is physically impossible and reflects the absence of any conservation constraint; for Tao $a_w$ 0.25 the first 38\,h of data contain no information about the later slow-down, so both mechanistic models place complete conversion earlier than observed and the neural ODE extrapolates in the opposite direction, which marks the information limit of a half-profile.
\inctab{fcst_main}
\begin{figure}[!htbp]\centering\incfig{fig11_forecast_bars}\caption{Forward extrapolation (training on $t\le0.5\,t_\mathrm{max}$); layout as in Figure~\ref{fig:loto}.}\label{fig:fcstbars}\end{figure}
\begin{figure}[!htbp]\centering\incfig{fig12_forecast}\caption{Forward extrapolation for three profiles; the shaded region was used for training. Lines connect the predictions at the held-out sampling times of the original method (dashed), PPK-FIP (solid) and the black-box neural ODE (dotted).}\label{fig:fcst}\end{figure}

\storylink{Leave-one-time-point-out validation is the decisive test of the paper, and PPK-FIP passes it: a \LotoAllRedOrig\,\% lower prediction error than the original method, significant in four cases and never significantly worse, with the gain concentrated where Step~2 found a missing mechanism. Forward extrapolation, with one split per case, places the two mechanistic estimators on a par and adds two pieces of information: the additional mechanisms need a profile that extends into the slow-down to be identified, and a half-profile without that information (Tao $a_w$ 0.25) bounds every estimator alike. It is also the protocol that separates mechanistic from black-box models in Step~6.}

\subsection{Step 4: Processivity and the Prior Carry the Gain; Deactivation and Averaging Do Not}
Table~\ref{tab:ablation} and Figure~\ref{fig:ablation} remove one component of PPK-FIP at a time, re-estimating every variant in every LOTO fold. Figure~\ref{fig:ablation}a shows the pooled LOTO RMSE of each variant against the dashed line of full PPK-FIP (\LotoAllNew), with the two original two-step forms in orange; Figure~\ref{fig:ablation}b shows the per-case ratio of each variant to PPK-FIP as a heat map, red where the variant is worse. Four results follow. (1)~Removing pseudo-processivity raises the pooled RMSE to \AblNoProc{} ($p$\,\AblNoProcP), and the heat map shows that the loss is almost entirely the Bornadel column (\AblNoProcBorn{} vs.\ \LotoBornNew, ratio 1.63), with small losses for Åkerman 100\,°C and Tao; this mechanism carries most of the out-of-sample gain. (2)~Within-batch deactivation is specific to the Åkerman 100\,°C profile: removing it leaves the pooled LOTO RMSE unchanged (\AblNoKd{} vs.\ \LotoAllNew) and gives the best pooled extrapolation (11.41 vs.\ 11.87), chiefly through the Mao batches, whose half-profiles do not identify $k_d$. The shrinkage prior therefore does what it was designed to do: it keeps $k_d$ at the prior level wherever a profile shows no late slow-down (Table~\ref{tab:params}), so that PPK-FIP with and without the term predict alike, and the term is available for the profiles, such as 100\,°C, that call for it. Removing both mechanisms gives \AblNoBoth, i.e.\ essentially the processivity effect. (3)~Removing the fingerprint prior (joint estimation with the same model averaging) gives \AblJoint, and with BIC selection instead of averaging 2.232; the heat map confines the loss to the sparse Åkerman column (ratios 1.24 and 1.53 at 70\,°C, 1.11 and 1.16 at 100\,°C), where nine or seven sampling times cannot pin down selectivities and processivity together. Conversely, the tight prior $\lambda=1$, which is closest to the two-step method, gives \AblLamOne{} ($p<0.001$) with its loss concentrated in Bornadel (ratio 1.12--1.13), because a tight prior holds the selectivities near their apparent fingerprint values. With $\lambda=3$ the result is insensitive to the processivity shrinkage $s_p$ (fixed $s_p=0.05$ or $0.2$ for all cases both give the meta-selected pooled RMSE to three decimals), and meta-validation chose $\lambda=3$ for every case. (4)~Model averaging over the simplified and full sets of lumped constants is not essential: the full set alone performs similarly (\AblFullOnly), whereas the simplified set alone is worse (\AblSimplOnly, ratio 1.32 for Tao $a_w$ 0.25). The two orange rows of the heat map, finally, show the original two-step forms red in almost every cell, with ratios of 1.73 for Bornadel and 1.55 for Tao $a_w$ 0.25.
\inctab{ablation}
\begin{figure}[!htbp]\centering\incfig{fig14_ablation}\caption{Ablation (LOTO). (a) Pooled RMSE of each variant; dashed line: full PPK-FIP. (b) Per-case RMSE ratio to PPK-FIP (red: worse than PPK-FIP; blue: better).}\label{fig:ablation}\end{figure}

A multi-task extension tests the opposite direction, adding rather than removing structure (Table~\ref{tab:pooling}, Figure~\ref{fig:pooling}). Profiles of the same reaction system share the composition path (L1), additionally the lumped constants with hierarchical shrinkage (L1h), or everything (L2), in the spirit of the original fingerprint step, which already pooled these profiles. Full pooling strongly improves Åkerman 100\,°C (\LotoAkHundredNew$\to$\PoolAkHundredLii) and both Tao cases (1.36$\to$1.26, 1.50$\to$1.30) but degrades Åkerman 70\,°C (1.89$\to$2.43), because the two temperatures share a path but not the deactivation; with the pooling level chosen by leave-one-case-out meta-validation the pooled RMSE is \PoolAllSel{} versus \PoolAllLo{} without pooling, i.e.\ no net gain, and pooling was not adopted.
\inctab{pooling}
\begin{figure}[!htbp]\centering\incfig[0.85\linewidth]{fig17_pooling}\caption{Multi-task pooling across conditions of the same system (LOTO RMSE).}\label{fig:pooling}\end{figure}

\storylink{The ablation turns a comparison into an attribution. Of the four components of PPK-FIP, pseudo-processivity and the calibrated prior are necessary; deactivation is a profile-specific option that the prior switches off where it is not needed; and model averaging can be replaced by the full lumped-constant set without loss. The dose--response along the prior width, from $\lambda=1$ (2.21) through $\lambda=3$ (2.12) to a flat prior (2.18), is the empirical answer to the question posed in the Introduction: the fingerprint should be used as information of calibrated width, neither as a constraint nor discarded. Pooling shows that sharing the path across conditions, the fingerprint idea in another guise, pays where the conditions differ only in time-scale (Tao), whereas per-profile estimation is preferable where they also differ in mechanism (Åkerman).}

\subsection{Step 5: The Added Mechanisms Behave Like Properties of the Reaction Systems}
If processivity is a property of the reaction system rather than a fitting device, a composition path learned at some conditions should transfer to others. Table~\ref{tab:transfer} and Figure~\ref{fig:transfer}a report the RMSE of each held-out profile predicted against $c$ from the other profiles of the same system, with the original path ($p_M=p_D=0$) and with the processive path. The processive path was better for all three Tao water activities (by 11--31\,\%), equal for Åkerman 70\,°C and the three Mao batches, and for Åkerman 100\,°C (2.26 vs.\ 2.63) it shows the boundary of what a $c$-domain path can carry; overall the RMSE ratio was \TransferRatio. Figure~\ref{fig:transfer}b shows the Tao $a_w$ 0.45 profile predicted from $a_w$ 0.25 and 0.35: the original path over-predicts M below $c=100$ and under-predicts D, the same signature as in Bornadel on a smaller scale, and the processive path removes it. The Tao processivity is small ($p_D\approx0.05$--$0.06$) but consistent across the three water activities (Table~\ref{tab:params}), which is why it transfers; the two Åkerman temperatures differ in deactivation, a time-domain property that the $c$-domain path, by construction, does not carry and that PPK-FIP estimates per profile instead.
\inctab{transfer}
\begin{figure}[!htbp]\centering\incfig{fig16_transfer}\caption{(a) Cross-condition transfer of the composition path. (b) Example: the profile at $a_w=0.45$ predicted from the profiles at $a_w=0.25$ and $0.35$ (dashed: original path; solid: processive path).}\label{fig:transfer}\end{figure}

The Monte-Carlo experiment answers the question with known truth. Table~\ref{tab:sim} and Figure~\ref{fig:sim} summarise \NSimReps{} synthetic replicates per scenario; the upper row of Figure~\ref{fig:sim} shows the RMSE of the fitted curves against the noise-free truth and the lower row the log-bias of $\hat S_M$, for six estimators and three truths. When the truth is the original model (S1), PPK-FIP and the original method recover the true curves equally well (\SimSiNew{} vs.\ \SimSiOrig{} molar\,\%, within one standard deviation), PPK-FIP's $S_M$ is unbiased ($+0.025$ against $+0.067$ for the two-step method), and the unregularised joint fit over-fits (\SimSiJoint, with an $S_M$ spread four times wider): the fingerprint-informed prior is what prevents spurious complexity. When the truth is pseudo-processive (S2), the original method is biased in a way no amount of data can cure: $\hat S_M$ is too high by \SimSMBiasOrigPct\,\% ($\ln$-bias \SimSMBiasOrig) with a curve error of \SimSiiOrig, the tight prior $\lambda=1$ retains half of that bias ($+0.22$), and PPK-FIP recovers the truth with a bias of \SimSMBiasNew{} and a curve error of \SimSiiNew. With within-batch deactivation (S3) the curve error falls from \SimSiiiOrig{} (original) and \SimSiiiFull{} (original full) to \SimSiiiNew; the full $F$ set can mimic a slow-down but only at the price of a biased $S_D$ ($-0.19$), while the variants with an explicit $k_d$ are unbiased.
\inctab{sim}
\begin{figure}[!htbp]\centering\incfig{fig15_simulation}\caption{Monte-Carlo experiment: curve RMSE against the noise-free truth (top) and log-bias of $\hat S_M$ (bottom) for three truths. PPK-FIP uses the meta-selected $\lambda=3$, $s_p=0.2$; FIP with $\lambda=1$ shown for comparison.}\label{fig:sim}\end{figure}

\storylink{Step~5 closes the loop on the ``apparent selectivities'' of the Introduction. With known truth, the two-step estimator overstates $S_M$ by a factor of 1.8 when pseudo-processivity is present, the same order as the shift from 0.909 to \BornSM{} found in the Bornadel data, and PPK-FIP recovers it; when the truth is the original model, PPK-FIP neither invents processivity nor loses accuracy. A generalisation of a published model must pass both tests. The transfer experiment adds real-data support that the processive path is a system property in the Tao case, and delineates the scope of the $c$-domain idea: it carries what is independent of the time-scale, while time-domain mechanisms such as deactivation are estimated per profile.}

\subsection{Step 6: Mechanism Is Still Necessary, and Machine Learning Is Most Useful on Top of It}
Tables~\ref{tab:lotobb} and \ref{tab:fcstbb} and Figure~\ref{fig:blackbox} compare the mechanistic estimators with four data-driven models given the same information (time, condition descriptor, initial state). In leave-one-time-point-out validation (Figure~\ref{fig:blackbox}a, log axis) the picture splits by sampling density. On the sparse Åkerman profiles all black-box models fail, with RMSE 13--21 for the Gaussian process, the MLP ensemble and TabICL and 3--4 for the neural ODE, against 1.9--3.5 for the mechanistic models. On the densely sampled profiles the Gaussian process (Bornadel 0.68, Tao $a_w$ 0.35/0.45 0.74, Mao 1.35) and the neural ODE (1.08, 1.46, 1.53) interpolate as well as or better than any mechanistic model, and because the 141 Mao observations dominate the pooled count, the neural ODE's pooled RMSE (\LotoBBNODE) lies below that of PPK-FIP (\LotoAllNew) even though PPK-FIP is lower on every sparse profile. The Mao result is the one anticipated in Step~2: a model that is free to leave mass conservation can follow the batch-1 plateau, whereas a mechanistic model reproduces the kinetics and leaves that feature to the measurement. In forward extrapolation (Figure~\ref{fig:blackbox}b) every black-box model fails in every case, with per-case errors of 6--47 and pooled RMSE of \FcstBBGP{} (GP), \FcstBBMLP{} (MLP ensemble), \FcstBBTabICL{} (TabICL) and \FcstBBNODE{} (neural ODE) versus \FcstAllNew{} for PPK-FIP, because nothing in them encodes the conservation of TMP or the sequential Ping Pong structure.
\inctab{loto_blackbox}\inctab{fcst_blackbox}
\begin{figure}[!htbp]\centering\incfig{fig13_blackbox}\caption{Mechanistic models versus black-box baselines (log scale): Gaussian process, MLP ensemble, pretrained TabICL foundation model and neural ODE.}\label{fig:blackbox}\end{figure}

The hybrid PPK-FIP-GP adds a Gaussian-process discrepancy to the mechanistic mean (Table~\ref{tab:hybrid}, Figure~\ref{fig:hybrid}). In leave-one-time-point-out validation (Figure~\ref{fig:hybrid}a) it has the lowest error of all methods considered, \HybLotoAll{} molar\,\%, i.e.\ \HybLotoRedOrig\,\% below the original method (\LotoAllOrig, $p$\,\HybLotoPOrig; lower fold error in \FoldHybWinsOrig{} of \FoldN{} folds), below PPK-FIP itself (\HybLotoBase) and below the best black-box model, the neural ODE (\LotoBBNODE); the gain is largest for the dense profiles (Bornadel 1.05$\to$0.73, Mao 3.10$\to$\HybLotoMao) and negligible for the sparse ones. In extrapolation (Figure~\ref{fig:hybrid}b) the PPK-FIP-GP bars coincide with the PPK-FIP bars (\HybFcstAll{} vs.\ \HybFcstBase), because the GP reverts to zero away from the data and the mechanistic model takes over, whereas the GP alone (pink) is off by a factor of two to thirty. Two further observations sharpen the interpretation. The same discrepancy added to the original two-step model also lowers its interpolation error (\HybLotoOrigGP, $p=\HybLotoPOrigGP$ against PPK-FIP-GP), which shows that the GP term is a general-purpose correction for digitisation and systematic error that any reasonable mechanistic mean can carry. What distinguishes the two hybrids is extrapolation: PPK-FIP-GP inherits the mechanistic gains of PPK-FIP (Åkerman 100\,°C 5.5 vs.\ 10.8 for ORIG+GP; pooled \HybFcstAll{} vs.\ 12.02), so the mechanistic improvements are what count outside the data window.
\inctab{hybrid}
\begin{figure}[!htbp]\centering\incfig{fig18_hybrid}\caption{Mechanistic models with and without a Gaussian-process discrepancy term, and the GP alone, in (a) leave-one-time-point-out validation and (b) forward extrapolation (log scale).}\label{fig:hybrid}\end{figure}

\storylink{Step~6 answers the question left open by Step~3, whether PPK-FIP's gain is just flexibility, by separating two kinds of accuracy. Inside the experimental window a discrepancy term can patch any reasonable mechanistic mean, and interpolation accuracy says little about the mechanism; outside the window, and for the selectivities themselves, only the mechanism counts, and there black-box learners fail while PPK-FIP keeps its advantage. The practical recommendation follows: PPK-FIP for parameter estimation and extrapolation, PPK-FIP-GP when the goal is the most accurate prediction within the experimental time window.}
